% ---------------------------------------------------------------------------
\section{Related work}
\label{sec:related}

The closest precedents to a bearing-faithful summarising lens are mostly not
called lenses. They sit in five literatures: interactive lenses and their design
spaces; labelling around a focus region; off-screen indicators on the viewport
border; radial and necklace cartography; and glyph maps with their underlying
local statistics.

\subsection{Interactive lenses and their design spaces}

%%RELATED-LENSES%%

\subsection{Marks around a focus region}

\emph{Excentric labelling} shows the labels of objects inside a cursor-centred
circle, connected to them by lines and updated as the cursor
moves~\cite{fekete1999excentric}. Labels are laid out in left and right stacks;
in its radial variant bearing is used only to order them. When there are too
many objects it shows a count and a sample of labels, and a glyph ``summarizing
the contents of the area'' is proposed as future work~\cite{fekete1999excentric}.
\emph{Extended excentric labelling} adds support for high and uneven density and
for summary statistics~\cite{bertini2009extended}
\doubt{what those statistics are and how they are drawn is unverified: full text
not accessed}.

The algorithmic literature on \emph{focus-region} and \emph{contour} labelling
places labels on the boundary of a circular focus
itself~\cite{fink2012focus,bekos2019external}. Fink et al.\ connect each label
to its site by a straight or Bézier leader; in their radial model a label's port
is the radial projection of its site onto the circle, and sites too close in
angle are resolved by keeping a maximum conflict-free subset; a facility-location
variant clusters the sites and labels one representative per cluster; and the
focus region's position can itself be optimised~\cite{fink2012focus}. Haunert and
Hermes require every leader to point at the centre, allow one position per
label, and solve the resulting maximum-weight independent set in real time in a
browser~\cite{haunert2014labeling}. Boundary labelling has been extended to
moving focus regions~\cite{heinsohn2014boundary}, to fisheye displacement of labels
out of a focus with optimisation once it stops~\cite{niedermann2019focus}, and to
labels drawn as arcs in an annular orbit around a circle with minimum total
leader length~\cite{bonerath2024orbital}, where a study with 54 participants found
straight leaders faster than orbital-radial ones at similar
accuracy~\cite{wallinger2026clarity}.
\doubt{Wallinger et al.'s numbers were read from arXiv v1, not the PacificVis
version of record.}

This is the geometric twin of what we do, and we build on it rather than around
it. The difference is in what the marks are and how conflicts are resolved. In
every one of these techniques a label stands for one object, a leader is always
drawn, and a conflict with the bearing is resolved either by \emph{selection}
(dropping a label) or by freeing the label's position. A summarising lens
places \emph{aggregates}, keeps every one of them, displaces them
\emph{minimally} from their bearing, and draws a leader only where the
resulting gap is large enough to mislead---so the leader's length is the
residual. We did not find a labelling technique that draws leaders conditionally
on displacement.

%%RELATED-OFFSCREEN%%

\subsection{Radial and necklace cartography}

%%RELATED-RADIAL%%

\subsection{Glyph maps, within-cell structure and local statistics}

Glyph placement strategies divide into data-driven and
structure-driven~\cite{ward2002taxonomy}; necklace placement is structure-driven
layout with a data-driven preferred position. Glyph design guidance is well
surveyed~\cite{borgo2013glyph}, but controlled evidence on glyph layout and on
glyphs over maps is thin~\cite{fuchs2017glyph}. Glyph-maps place small temporal
glyphs at the data's own grid locations~\cite{wickham2012glyphmaps};
\emph{tilemaps} and \emph{gridded glyphmaps} aggregate into a regular grid and
draw a multivariate glyph per cell~\cite{slingsby2018tilemaps,slingsby2023gridded,
laksono2024gridded}, and multivariate glyph placement with level of detail uses
hierarchical centroids~\cite{mcnabb2019multivariate}. Tilemaps were proposed as a
way to summarise the outputs of geographically weighted statistics, with the
grid interactively resized and panned to expose the modifiable areal unit
problem and a distance-decay kernel larger than the tiles as an
option~\cite{slingsby2018tilemaps}. Gridded glyphmaps keep the grid fixed in screen
space and re-aggregate on zoom, and name glyphs that convey within-cell
heterogeneity as future work~\cite{slingsby2023gridded}.

That heterogeneity is what two hexagon techniques encode. Honeycomb Plots
observe that colour-coded hexbins cannot distinguish ``uniform distributions
within tiles \ldots from clusters or trends if their numbers of points match'',
and cut each tile along the regression plane of its within-tile density; a
between-subjects study with 42 participants found strong evidence for slope
tasks~\cite{trautner2022honeycomb}. HexTiles derive a per-tile confidence from
weighted within-tile variance, argued as mitigation of the modifiable areal unit
problem~\cite{kawakami2024hextiles}; its weighting is described only in words,
and its user study (seven participants) tested the tiles without the variance
encoding. Polar histograms of street bearings summarise a city's orientation
order~\cite{boeing2019urban}; they bin the bearing \emph{of} edges rather than
bearing \emph{from} a chosen centre, which is the quantity a lens bins.

%%RELATED-STATS%%
